%% kapa: exam=1 round=31 year=2020 period=2 form=A subject=appraisal_law subject_ko=감정평가관계법규 seq=22 number=22
%% source: Q-Net 감정평가사 기출문제 (cst00309) artlSeq=5207219 file=extracted/1008/5207219/A형/A형/2020년도 제31회 감정평가사 1차 시험_2교시_A형.pdf
%% conversion: kordoc 4.24.51 -> tools/gen1.py (auto); review=proofread-sonnet; flags=none
\begin{question}{22}
국유재산법령상 일반재산에 관한 설명으로 옳은 것은?
\begin{choices}
  \item 총괄청은 일반재산의 관리ㆍ처분에 관한 사무의 일부를 위탁받을 수 없다.
  \item 증권을 제외한 일반재산을 지방자치단체에 처분할 때 처분재산의 예정가격은 두 개의 감정평가업자의 평가액을 산술평균한 금액으로 결정한다.
  \item 조림을 목적으로 하는 토지의 대부기간은 25년 이상으로 한다.
  \item 중앙관서의 장은 일반재산을 교환하려면 그 내용을 감사원에 보고하여야 한다.
  \item 일반재산을 매각하면서 매각대금을 한꺼번에 납부하기로 한 경우 매각대금의 완납 이전에도 해당 매각재산의 소유권 이전이 가능하다.
\end{choices}
\end{question}
\answer{4}
